% !TEX spellcheck = en_US
% !TEX spellcheck = LaTeX
\documentclass[letterpaper,10pt,english]{article}
%\usepackage{%
amsfonts,%
amsmath,%
amssymb,%
amsthm,%
babel,%
bbm,%
%biblatex,%
caption,%
centernot,%
color,%
enumerate,%
%enumitem,%
epsfig,%
epstopdf,%
etex,%
fancybox,%
framed,%
fullpage,%
%geometry,%
graphicx,%
hyperref,%
latexsym,%
mathptmx,%
mathtools,%
multicol,%
paralist,%
pgf,%
pgfplots,%
pgfplotstable,%
pgfpages,%
proof,%
psfrag,%
%subfigure,%
tcolorbox,%
tfrupee,%
tikz,%
times,%
%ulem,%
url,%
xcolor,%
mathpazo,%
}

\definecolor{shadecolor}{gray}{.95}%{rgb}{1,0,0}
\usepackage[margin=1in,top=0.75in]{geometry}
\usepackage[mathscr]{eucal}
\usepgflibrary{shapes}
\usepgfplotslibrary{fillbetween}
\usetikzlibrary{%
arrows,%
backgrounds,%
chains,%
decorations.pathmorphing,% /pgf/decoration/random steps | erste Graphik
decorations.text,% 
matrix,%
positioning,% wg. " of "
fit,%
patterns,%
petri,%
plotmarks,%
scopes,%
shadows,%
shapes.misc,% wg. rounded rectangle
shapes.arrows,%
shapes.callouts,%
shapes%
}

%\pgfplotsset{compat=newest} %<------ Here
\pgfplotsset{compat=1.11} %<------ Or use this one

\theoremstyle{plain}
\newtheorem{thm}{Theorem}[section]
\newtheorem{lem}[thm]{Lemma}
\newtheorem{prop}[thm]{Proposition}
\newtheorem{cor}[thm]{Corollary}
\newtheorem{clm}[thm]{Claim}

\theoremstyle{definition}
\newtheorem{defn}[thm]{Definition}
\newtheorem{conj}[thm]{Conjecture}
\newtheorem{exmp}[thm]{Example}
\newtheorem{axiom}[thm]{Axiom}
\newtheorem{assum}[thm]{Assumption}
\newtheorem{exerc}[thm]{Exercise}
\newtheorem*{defin}{Definition}

\theoremstyle{remark}
\newtheorem{rem}{Remark}
\newtheorem{note}{Note}

\newcommand{\iid}{\emph{i.i.d.}\ }
\newcommand{\Cov}{\operatorname{Cov}}
%\newcommand{\det}{\operatorname{det}}
%\newcommand{\Real}{\mathbb{R}}
\newcommand{\tr}{\operatorname{tr}}
\newcommand{\Var}{\operatorname{Var}}
\newcommand{\cov}{\operatorname{cov}}

%\renewcommand{\proof}[1]{\begin{proof}#1\end{proof}}
\newcommand{\EQ}[1]{\begin{equation*}#1\end{equation*}}
\newcommand{\EQN}[1]{\begin{equation}#1\end{equation}}
\newcommand{\eq}[1]{\begin{align*}#1\end{align*}}
\newcommand{\meq}[2]{\begin{xalignat*}{#1}#2\end{xalignat*}}
\newcommand{\norm}[1]{\left\lVert#1\right\rVert}
\newcommand{\inner}[1]{\left\langle #1\right\rangle}
\newcommand{\abs}[1]{\left\lvert#1\right\rvert}
\newcommand{\expect}[1]{\mathbb{E}\left[{#1}\right]}
\newcommand{\prob}[1]{\mathbb{P}\left[{#1}\right]}
\newcommand{\given}{\; \big\vert \;} 
\newcommand{\set}[1]{\left\{#1\right\}} 
\newcommand{\indicator}[1]{1_{\set{#1}}} 
\newcommand{\Ind}[1]{\mathbbm{1}_{#1}} 
\newcommand{\SetIn}[1]{\mathbbm{1}_{\set{#1}}} 
\newcommand{\red}[1]{\textcolor{red}{#1}} 
\newcommand{\floor}[1]{\left\lfloor#1\right\rfloor}
\newcommand{\ceil}[1]{\left\lceil#1\right\rceil}


\newcommand{\C}{\mathbb{C}}
\newcommand{\D}{\mathbb{D}}
\newcommand{\E}{\mathbb{E}}
\newcommand{\F}{\mathbb{F}}
\newcommand{\N}{\mathbb{N}}
\renewcommand{\P}{\mathbb{P}}
\newcommand{\Q}{\mathbb{Q}}
\newcommand{\R}{\mathbb{R}}
\newcommand{\Z}{\mathbb{Z}}

\newcommand{\bA}{\mathbf{A}}
\newcommand{\bI}{\mathbf{I}}
\newcommand{\bU}{\mathbf{1}}
\newcommand{\bx}{\mathbf{x}}
\newcommand{\bX}{\mathbf{X}}
\newcommand{\bZ}{\mathbf{Z}}


\newcommand{\cA}{\mathcal{A}}
\newcommand{\cB}{\mathcal{B}}
\newcommand{\cC}{\mathcal{C}}
\newcommand{\cF}{\mathcal{F}}
\newcommand{\cG}{\mathcal{G}}
\newcommand{\cM}{\mathcal{M}}
\newcommand{\cN}{\mathcal{N}}
\newcommand{\cP}{\mathcal{P}}
\newcommand{\cT}{\mathcal{T}}
\newcommand{\cX}{\mathcal{X}}
\newcommand{\cY}{\mathcal{Y}}

\newcommand{\sA}{\mathscr{A}}
\newcommand{\sB}{\mathscr{B}}
\newcommand{\sC}{\mathscr{C}}
\newcommand{\sE}{\mathscr{E}}
\newcommand{\sF}{\mathscr{F}}
\newcommand{\sG}{\mathscr{G}}
\newcommand{\sH}{\mathscr{H}}
\newcommand{\sL}{\mathscr{L}}
\newcommand{\sS}{\mathscr{S}}
\newcommand{\sT}{\mathscr{T}}
\newcommand{\sX}{\mathscr{X}}
\newcommand{\sY}{\mathscr{Y}}
\newcommand{\sZ}{\mathscr{Z}}

\newcommand{\tS}{\tilde{S}}
% Debug
\newcommand{\todo}[1]{\begin{color}{blue}{{\bf~[TODO:~#1]}}\end{color}}

% a few handy macros

\renewcommand{\le}{\leqslant}
\renewcommand{\ge}{\geqslant}
\newcommand\matlab{{\sc matlab}}
\newcommand{\goto}{\rightarrow}
\newcommand{\bigo}{{\mathcal O}}
%\newcommand{\half}{\frac{1}{2}}
%\newcommand\implies{\quad\Longrightarrow\quad}
\newcommand\reals{{{\rm l} \kern -.15em {\rm R} }}
\newcommand\complex{{\raisebox{.043ex}{\rule{0.07em}{1.56ex}} \hskip -.35em {\rm C}}}


% macros for matrices/vectors:

% matrix environment for vectors or matrices where elements are centered
\newenvironment{mat}{\left[\begin{array}{ccccccccccccccc}}{\end{array}\right]}
\newcommand\bcm{\begin{mat}}
\newcommand\ecm{\end{mat}}

% matrix environment for vectors or matrices where elements are right justifvied
\newenvironment{rmat}{\left[\begin{array}{rrrrrrrrrrrrr}}{\end{array}\right]}
\newcommand\brm{\begin{rmat}}
\newcommand\erm{\end{rmat}}

% for left brace and a set of choices
%\newenvironment{choices}{\left\{ \begin{array}{ll}}{\end{array}\right.}
\newcommand\when{&\text{if~}}
\newcommand\otherwise{&\text{otherwise}}
% sample usage:
%\delta_{ij} = \begin{choices} 1 \when i=j, \\ 0 \otherwise \end{choices}


% for labeling and referencing equations:
\newcommand{\eql}{\begin{equation}\label}
\newcommand{\eqn}[1]{(\ref{#1})}
% can then do
%\eql{eqnlabel}
%...
%\end{equation}
% and refer to it as equation \eqn{eqnlabel}.
% some useful macros for finite difference methods:
\newcommand\unp{U^{n+1}}
\newcommand\unm{U^{n-1}}
% for chemical reactions:
\newcommand{\react}[1]{\stackrel{K_{#1}}{\rightarrow}}
\newcommand{\reactb}[2]{\stackrel{K_{#1}}{~\stackrel{\rightleftharpoons}
 {\scriptstyle K_{#2}}}~}
\makeatletter
\def\th@plain{%
\thm@notefont{}% same as heading font
\itshape % body font
}
\def\th@definition{%
\thm@notefont{}% same as heading font
\normalfont % body font
}
\makeatother
\date{}



\usepackage{%
	amsfonts,%
	amsmath,%	
	amssymb,%
	amsthm,%
	babel,%
	bbm,%
	%biblatex,%
	caption,%
	centernot,%
	color,%
	enumerate,%
	%enumitem,%
	epsfig,%
	epstopdf,%
	etex,%
	fancybox,%
	framed,%
	fullpage,%
	%geometry,%
	graphicx,%
	hyperref,%
	latexsym,%
	mathptmx,%
	mathtools,%
	multicol,%
	paralist,%
	pgf,%
	pgfplots,%
	pgfplotstable,%
	pgfpages,%
	proof,%
	psfrag,%
	%subfigure,%	
	tcolorbox,%
	tfrupee,%
	tikz,%
	times,%
	%ulem,%
	url,%
	xcolor,%
	mathpazo,%
}

\definecolor{shadecolor}{gray}{.95}%{rgb}{1,0,0}
\usepackage[margin=1in,top=0.75in]{geometry}
\usepackage[mathscr]{eucal}
\usepgflibrary{shapes}
\usepgfplotslibrary{fillbetween}
\usetikzlibrary{%
  arrows,%
  backgrounds,%
  chains,%
  decorations.pathmorphing,% /pgf/decoration/random steps | erste Graphik
  decorations.text,% 
  matrix,%
  positioning,% wg. " of "
  fit,%
  patterns,%
  petri,%
  plotmarks,%
  scopes,%
  shadows,%
  shapes.misc,% wg. rounded rectangle
  shapes.arrows,%
  shapes.callouts,%
  shapes%
}

%\pgfplotsset{compat=newest} %<------ Here
\pgfplotsset{compat=1.11} %<------ Or use this one

\theoremstyle{plain}
\newtheorem{thm}{Theorem}[section]
\newtheorem{lem}[thm]{Lemma}
\newtheorem{prop}[thm]{Proposition}
\newtheorem{cor}[thm]{Corollary}
\newtheorem{clm}[thm]{Claim}

\theoremstyle{definition}
\newtheorem{defn}[thm]{Definition}
\newtheorem{conj}[thm]{Conjecture}
\newtheorem{exmp}[thm]{Example}
\newtheorem{axiom}[thm]{Axiom}
\newtheorem{assum}[thm]{Assumption}
\newtheorem{exerc}[thm]{Exercise}
\newtheorem*{defin}{Definition}

\theoremstyle{remark}
\newtheorem{rem}{Remark}
\newtheorem{note}{Note}

\newcommand{\iid}{\emph{i.i.d.}\ }
\newcommand{\Cov}{\operatorname{Cov}}
%\newcommand{\det}{\operatorname{det}}
%\newcommand{\Real}{\mathbb{R}}
\newcommand{\tr}{\operatorname{tr}}
\newcommand{\Var}{\operatorname{Var}}
\newcommand{\cov}{\operatorname{cov}}

%\renewcommand{\proof}[1]{\begin{proof}#1\end{proof}}
\newcommand{\EQ}[1]{\begin{equation*}#1\end{equation*}}
\newcommand{\EQN}[1]{\begin{equation}#1\end{equation}}
\newcommand{\eq}[1]{\begin{align*}#1\end{align*}}
\newcommand{\meq}[2]{\begin{xalignat*}{#1}#2\end{xalignat*}}
\newcommand{\norm}[1]{\left\lVert#1\right\rVert}
\newcommand{\inner}[1]{\left\langle #1\right\rangle}
\newcommand{\abs}[1]{\left\lvert#1\right\rvert}
\newcommand{\expect}[1]{\mathbb{E}\left[{#1}\right]}
\newcommand{\prob}[1]{\mathbb{P}\left[{#1}\right]}
\newcommand{\given}{\; \big\vert \;} 
\newcommand{\set}[1]{\left\{#1\right\}} 
\newcommand{\indicator}[1]{1_{\set{#1}}} 
\newcommand{\Ind}[1]{\mathbbm{1}_{#1}} 
\newcommand{\SetIn}[1]{\mathbbm{1}_{\set{#1}}} 
\newcommand{\red}[1]{\textcolor{red}{#1}} 
\newcommand{\floor}[1]{\left\lfloor#1\right\rfloor}
\newcommand{\ceil}[1]{\left\lceil#1\right\rceil}


\newcommand{\C}{\mathbb{C}}
\newcommand{\D}{\mathbb{D}}
\newcommand{\E}{\mathbb{E}}
\newcommand{\F}{\mathbb{F}}
\newcommand{\N}{\mathbb{N}}
\renewcommand{\P}{\mathbb{P}}
\newcommand{\Q}{\mathbb{Q}}
\newcommand{\R}{\mathbb{R}}
\newcommand{\Z}{\mathbb{Z}}

\newcommand{\bA}{\mathbf{A}}
\newcommand{\bI}{\mathbf{I}}
\newcommand{\bU}{\mathbf{1}}
\newcommand{\bx}{\mathbf{x}}
\newcommand{\bX}{\mathbf{X}}
\newcommand{\bZ}{\mathbf{Z}}


\newcommand{\cA}{\mathcal{A}}
\newcommand{\cB}{\mathcal{B}}
\newcommand{\cC}{\mathcal{C}}
\newcommand{\cF}{\mathcal{F}}
\newcommand{\cG}{\mathcal{G}}
\newcommand{\cM}{\mathcal{M}}
\newcommand{\cN}{\mathcal{N}}
\newcommand{\cP}{\mathcal{P}}
\newcommand{\cT}{\mathcal{T}}
\newcommand{\cX}{\mathcal{X}}
\newcommand{\cY}{\mathcal{Y}}

\newcommand{\sA}{\mathscr{A}}
\newcommand{\sB}{\mathscr{B}}
\newcommand{\sC}{\mathscr{C}}
\newcommand{\sE}{\mathscr{E}}
\newcommand{\sF}{\mathscr{F}}
\newcommand{\sG}{\mathscr{G}}
\newcommand{\sH}{\mathscr{H}}
\newcommand{\sL}{\mathscr{L}}
\newcommand{\sS}{\mathscr{S}}
\newcommand{\sT}{\mathscr{T}}
\newcommand{\sX}{\mathscr{X}}
\newcommand{\sY}{\mathscr{Y}}
\newcommand{\sZ}{\mathscr{Z}}

\newcommand{\tS}{\tilde{S}}
% Debug
\newcommand{\todo}[1]{\begin{color}{blue}{{\bf~[TODO:~#1]}}\end{color}}

% a few handy macros

\renewcommand{\le}{\leqslant}
\renewcommand{\ge}{\geqslant}
\newcommand\matlab{{\sc matlab}}
\newcommand{\goto}{\rightarrow}
\newcommand{\bigo}{{\mathcal O}}
%\newcommand{\half}{\frac{1}{2}}
%\newcommand\implies{\quad\Longrightarrow\quad}
\newcommand\reals{{{\rm l} \kern -.15em {\rm R} }}
\newcommand\complex{{\raisebox{.043ex}{\rule{0.07em}{1.56ex}} \hskip -.35em {\rm C}}}


% macros for matrices/vectors:

% matrix environment for vectors or matrices where elements are centered
\newenvironment{mat}{\left[\begin{array}{ccccccccccccccc}}{\end{array}\right]}
\newcommand\bcm{\begin{mat}}
\newcommand\ecm{\end{mat}}

% matrix environment for vectors or matrices where elements are right justifvied
\newenvironment{rmat}{\left[\begin{array}{rrrrrrrrrrrrr}}{\end{array}\right]}
\newcommand\brm{\begin{rmat}}
\newcommand\erm{\end{rmat}}

% for left brace and a set of choices
%\newenvironment{choices}{\left\{ \begin{array}{ll}}{\end{array}\right.}
\newcommand\when{&\text{if~}}
\newcommand\otherwise{&\text{otherwise}}
% sample usage:
%  \delta_{ij} = \begin{choices} 1 \when i=j, \\ 0 \otherwise \end{choices}


% for labeling and referencing equations:
\newcommand{\eql}{\begin{equation}\label}
\newcommand{\eqn}[1]{(\ref{#1})}
% can then do
%  \eql{eqnlabel}
%  ...
%  \end{equation}
% and refer to it as equation \eqn{eqnlabel}.  
% some useful macros for finite difference methods:
\newcommand\unp{U^{n+1}}
\newcommand\unm{U^{n-1}}
% for chemical reactions:
\newcommand{\react}[1]{\stackrel{K_{#1}}{\rightarrow}}
\newcommand{\reactb}[2]{\stackrel{K_{#1}}{~\stackrel{\rightleftharpoons}
   {\scriptstyle K_{#2}}}~}
\makeatletter
\def\th@plain{%
  \thm@notefont{}% same as heading font
  \itshape % body font
}
\def\th@definition{%
  \thm@notefont{}% same as heading font
  \normalfont % body font
}
\makeatother
\date{}





\title{Lecture-03: Independence}
\author{}

\begin{document}
\maketitle

\section{Law of Total Probability}
\begin{tcolorbox}
\begin{exerc}[Countably infinite coin tosses] 
\label{exerc:CountableCoinToss}
Consider a sequence of coin tosses, such that the sample space is $\Omega = \set{H,T}^\N$. 
For set of outcomes $E_n \triangleq \set{\omega \in \Omega: \omega_n = H}$, 
we consider an event space generated by $\sF \triangleq \sigma(\set{E_n: n\in \N})$. 
Let $\sF_n$ be the event space generate by the first $n$ coin tosses, 
i.e. $\sF_n \triangleq \sigma(\set{E_i: i \in [n]})$.  
Let $A_n$ be the set of outcomes corresponding to at least one head in first $n$ outcomes 
$
A_n \triangleq \set{\omega \in \Omega: \omega_i = H \text{ for some } i \in [n]} = \cup_{i=1}^nE_i \in \sF,
$
and $B_n$ be the set of outcomes corresponding to first head at the $n$th outcome
$
B_n \triangleq \set{\omega \in \Omega: \omega_1 = \dots = \omega_{n-1}= T, \omega = H} 
= \cap_{i=1}^{n-1}E_i^c\cap E_n \in \sF.
$
\begin{enumerate}
\item Show that $\sF = \sigma(\set{\sF_n: n \in \N})$. 
\item Show that $\sigma(\set{A_n: n \in \N}) \subseteq \sF$ and $\sigma(\set{B_n: n\in \N}) \subseteq \sF$. 
\end{enumerate}
\end{exerc}
\end{tcolorbox}
\begin{thm}[Law of total probability] 
For a probability space $(\Omega, \sF, P)$, 
consider a sequence of events $B\in \sF^\N$ that partitions the sample space $\Omega$, i.e. 
$B_m \cap B_n = \emptyset$ for all $m \neq n$, and $\cup_{n \in \N}B_n = \Omega$. 
Then, for any event $A \in \sF$, we have 
\EQ{
P(A) = \sum_{n \in \N}P(A \cap B_n).
}
\end{thm}
\begin{proof}
We can expand any event $A \in \sF$ in terms of any partition $B$ of the sample space $\Omega$ as 
\EQ{
A = A \cap \Omega = A \cap (\cup_{n \in \N}B_n) = \cup_{n \in \N}(A \cap B_n).
}
From the mutual disjointness of the events $B \in \sF^\N$, it follows that the sequence $(A \cap B_n \in \sF: n \in \N)$ is mutually disjoint. 
The result follows from the countable additivity of probability of disjoint events. 
\end{proof}

\begin{shaded}
\begin{exmp}[Countably infinite coin tosses] 
Consider the sample space $\Omega = \set{H,T}^\N$ and event space $\sF$ generated by sequence $E \in \sF^\N$ defined in Exercise~\ref{exerc:CountableCoinToss}. 
We observe that any event $A \in \sF_n$ can be written as 
\begin{equation*}
A = \cup_{\omega \in A}\set{\omega} 
= \cup_{\omega \in A}\cap_{i=1}^n(\set{\omega \in E_i}\cup\set{\omega \notin E_i}).
\end{equation*}  
\end{exmp}
\end{shaded}

\section{Independence}
\begin{defn}[Independence of events] 
For a probability space $(\Omega, \sF, P)$, a family of events $A \in \sF^I$ is said to be independent, 
if for any finite set $F \subseteq I$, we have 
\EQ{
P(\cap_{i \in F}A_i) = \prod_{i \in F}P(A_i). 
}
\end{defn}
\begin{rem} 
The certain event $\Omega$ and the impossible event $\emptyset$ are always independent to every event $A \in \sF$. 
\end{rem}
\begin{shaded}
\begin{exmp}[Two coin tosses] 
Consider two coin tosses, such that the sample space is $\Omega = \set{HH, HT, TH, TT}$, 
and the event space is $\sF = \cP(\Omega)$. 
It suffices to define a probability function $P: \sF \to [0,1]$ on the sample space. 
We define one such probability function $P$, such that 
\EQ{
P(\set{HH}) = P(\set{HT}) = P(\set{TH}) =P(\set{TT}) = \frac{1}{4}. 
}
Let event $E_1 \triangleq \set{HH, HT}$ and $E_2 \triangleq \set{HH, TH}$ correspond to getting a head on the first or the second toss respectively. 

From the defined probability function, we obtain the probability of getting a tail on the first or the second toss is $\frac{1}{2}$, 
and identical to the probability of getting a head on the first or the second toss. 
That is, $P(E_1) = P(E_2) = \frac{1}{2}$ and the intersecting event $E_1\cap E_2 = \set{HH}$ with the probability $P(E_1 \cap E_2) = \frac{1}{4}$. 
That is, for events $E_1, E_2 \in \sF$, we have 
\EQ{
P(E_1 \cap E_2) = P(E_1)P(E_2).
}
That is, events $E_1$ and $E_2$ are independent. 
\end{exmp}
\end{shaded}
\begin{shaded}
\begin{exmp}[Countably infinite coin tosses] 
Consider the outcome space $\Omega = \set{H,T}^\N$ and event space $\sF$ generated by the sequence $E$ defined in Exercise~\ref{exerc:CountableCoinToss}.  
We define a probability function $P:\sF \to [0,1]$ by $P(\cap_{i\in F}E_i)= p^{\abs{F}}$ for any finite subset $F \subseteq \N$.  
By definition, $E \in \sF^\N$ is a sequence of independent events. 
Consider $A, B \in \sF^\N$, 
where $A_n \triangleq \cup_{i=1}^nE_i$ and $B_n \triangleq \cap_{i=1}^{n-1}E_i^c\cap E_n \in \sF$ for all $n\in\N$.
It follows that $P(A_n) = 1-(1-p)^n$ and $P(B_n) = p(1-p)^{n-1}$ for $n \in \N$. 

For any $\omega \in \Omega$, we can define the number of heads in first $n$ trials by $k_n(\omega) \triangleq\sum_{i=1}^n\SetIn{H}(\omega_i) = \sum_{i=1}^n\SetIn{\omega \in E_i}$. 
For any general event $A \in \sF_n = \sigma(\set{E_i: i \in [n]})$, we can write 
\EQ{
P(A) = \sum_{\omega \in A} \prod_{i=1}^n\Big[P\set{\omega \in E_i}  + P\set{\omega \in E_i^c}\Big] 
= \sum_{\omega \in A}p^{k_n(\omega)}(1-p)^{n-k_n(\omega)}.
}  
\end{exmp}
\end{shaded}

%\begin{shaded}
%\begin{exmp}
%Consider a random experiments, 
%where a random point is thrown on non-negative real line $\R_+$. 
%%We denote the location of random point by $S$.   
%The outcome space of this experiment is $\R_+$ and the Borel event space is $\sB(\R_+)$. 
%Let 
%\end{exmp}
%\end{shaded}
\begin{shaded}
\begin{exmp}[Counter example] 
Consider a probability space $(\Omega, \sF, P)$ and the events $A_1, A_2, A_3 \in \sF$. 
The condition $P(A_1\cap A_2\cap A_3) = P(A_1)P(A_2)P(A_3)$ is not sufficient to guarantee independence of the three events. 
In particular, we see that if
\meq{2}{
&P(A_1\cap A_2\cap A_3) = P(A_1)P(A_2)P(A_3),&&P(A_1\cap A_2\cap A_3^c) \neq P(A_1)P(A_2)P(A_3^c), 
}
then $P(A_1\cap A_2) =  P(A_1\cap A_2\cap A_3) + P(A_1\cap A_2\cap A_3^c) \neq P(A_1)P(A_2)$.
\end{exmp} 
%\begin{proof}
%We will show this by contradiction.   
%Suppose $A_1, A_2 \in \sF$ are independent, then we can write by independence and the law of total probability 
%\EQ{
%P(A_1)P(A_2) = P(A_1\cap A_2) = P(A_1\cap A_2\cap A_3)  + P(A_1\cap A_2\cap A_3^c) 
%}
%Re-arranging the above terms and from the first hypothesis, we get 
%\EQ{
%P(A_1\cap A_2\cap A_3^c)  = P(A_1)P(A_2)(1 - P(A_3)) = P(A_1)P(A_2)P(A_3^c).
%}
%This contradicts the second hypothesis, and hence the events $A_1$ and $A_2$ can't be independent. 
%\end{proof}
\end{shaded}
\begin{defn}
A family of collections of events $(\sA_i \subseteq \sF: i \in I)$ is called independent,
if for any finite set $F \subseteq I$ and $A_i \in \sA_i$ for all $i \in F$, we have 
\EQ{
P(\cap_{i \in F}A_i) = \prod_{i \in F}P(A_i).
} 
\end{defn}

\section{Conditional Probability}
Consider $N$ trials of a random experiment over an outcome space $\Omega$ and an event space $\sF$.  
Let $\omega_n \in \Omega$ denote the outcome of the experiment of the $n$th trial. 
Consider two events $A, B \in \sF$ and denote the number of times event $A$ and event $B$ occurs by $N(A)$ and $N(B)$ respectively. 
We denote the number of times both events $A$ and $B$ occurred by $N(A\cap B)$. 
Then, we can write these numbers in terms of indicator functions as 
\meq{3}{
&N(A) = \sum_{n=1}^N\SetIn{\omega_n \in A}, &&N(B) = \sum_{n=1}^N\SetIn{\omega_n \in B}, &&N(A\cap B) = \sum_{n=1}^N\SetIn{\omega_n \in A \cap B}.
}
We denote the relative frequency of events $A, B, A\cap B$ in $N$ trials by $\frac{N(A)}{N}, \frac{N(B)}{N}, \frac{N(A\cap B)}{N}$ respectively.  
We can find the relative frequency of events $A$, on the trials where $B$ occurred as 
\EQ{
\frac{\frac{N(A\cap B)}{N}}{\frac{N(B)}{N}} = \frac{N(A \cap B)}{N(B)}.
} 
Inspired by the relative frequency, we define the conditional probability function conditioned on events. 
\begin{defn}
Fix an event $B \in \sF$ such that $P(B) > 0$, we can define the conditional probability $P(\cdot | B): \sF \to [0,1]$ of any event $A \in \sF$ conditioned on the event $B$ as 
\EQ{
P(A|B) = \frac{P(A\cap B)}{P(B)}.
}
\end{defn}
\begin{lem}[Conditional probability] 
For any event $B \in \sF$ such that $P(B) > 0$, the conditional probability $P(\cdot | B): \sF \to [0,1]$ is a probability measure on space $(\Omega, \sF)$. 
\end{lem}
\begin{proof} 
We will show that the conditional probability satisfies all three axioms of a probability measure.
\begin{enumerate}
\item  [\textbf{Non-negativity:}] For all events $A \in \sF$, we have $P(A|B) \ge 0$ since $P(A \cap B) \ge 0$. 
\item  [\textbf{$\sigma$-additivity:}] For an infinite sequence of mutually disjoint events $(A_i \in \sF: i \in \N)$ such that $A_i \cap A_j = \emptyset$ for all $i \neq j$, we have $P(\cup_{i \in \N}A_i|B) = \sum_{i \in \N}P(A_i|B)$. 
This follows from disjointness of the sequence $(A_i\cap B \in \sF: i \in \N)$. 
\item [\textbf{Certainty:}] Since $\Omega \cap B = B$, we have $P(\Omega |B) = 1$. 
\end{enumerate} 
\end{proof}
\begin{rem}
For two independent events $A, B \in \sF$ such that $P(A\cap B) > 0$, 
we have  $P(A|B) = P(A)$ and  $P(B|A) = P(B)$. 
%Therefore, if $P(B) > 0$ then, and if $P(A) > 0$ then $P(B|A) = P(B)$. 
If either $P(A) = 0$ or $P(B) = 0$, then $P(A \cap B) = 0$. 
\end{rem}
\begin{rem}
For any partition $B$ of the sample space $\Omega$, if $P(B_n) > 0$ for all $n \in \N$, then from the law of total probability and the definition of conditional probability, we have 
\EQ{
P(A) = \sum_{n \in \N}P(A|B_n)P(B_n).
}
\end{rem}


\section{Conditional Independence}
\begin{defn}[Conditional independence of events] 
For a probability space $(\Omega, \sF, P)$, a family of events $A \in \sF^I$ is said to be conditionally independent given an event $C \in \sF$ such that $P(C) > 0$, 
if for any finite set $F \subseteq I$, we have 
\EQ{
P(\cap_{i \in F}A_i|C) = \prod_{i \in F}P(A_i|C). 
}
\end{defn}
\begin{rem} 
Let $C \in \sF$ be an event such that $P(C) > 0$. 
Two events $A, B \in \sF$ are said to be conditionally independent given event $C$, if 
\EQ{
P(A \cap B|C) = P(A|C) P(B|C). 
}
If the event $C = \Omega$, it implies that $A, B$ are independent events. 
\end{rem}
\begin{rem}
Two events may be independent, but not conditionally independent and vice versa. 
\end{rem}
\begin{shaded}
\begin{exmp} 
Consider two independent events $A, B \in \sF$ such that $P(A\cap B) > 0$ and $P(A \cup B) < 1$. 
Then the events $A$ and $B$ are not conditionally independent given $A\cup B$.  
To see this, we observe that 
\EQ{
P(A \cap B | A\cup B) = \frac{P( (A \cap B)\cap (A \cup B))}{P(A\cup B)} = \frac{P(A \cap B)}{P(A\cup B)} 
= \frac{P(A)P(B)}{P(A \cup B)} =  {P(A|A \cup B)}{P(B)}. 
}
We further observe that $P(B|A\cup B) = \frac{P(B)}{P(A\cup B)} \neq P(B)$ and hence $P(A \cap B | A\cup B) \neq P(A|A\cup B)P(B|A\cup B)$. 
\end{exmp}
%\end{shaded}
%\begin{shaded}
\begin{exmp} 
Consider two non-independent events $A, B \in \sF$ such that $P(A) > 0$. 
Then the events $A$ and $B$ are conditionally independent given $A$. 
To see this, we observe that 
\EQ{
P(A \cap B | A) = \frac{P(A \cap B)}{P(A)} = P(B|A)P(A|A). 
}
\end{exmp}
\end{shaded}
\end{document}
